\section{Cloud und Datenschutz}

\subsection{Data Sovereignty}
Regionale Datenschutzvorgaben verlangen, dass Daten physisch in
bestimmten Gebieten gespeichert werden. Cloud-Anbieter nutzen zwar
geografisch verteilte Replikation für hohe Verfügbarkeit, jedoch kann
diese verteilte Speicherung gegen lokale Vorschriften verstoßen.
Eine Daten-Souveränitäts-Architektur stellt sicher, dass geschützte
Daten regelkonform abgelegt werden. Replikationsmechanismen müssen
daher Compliance-fähig konfigurierbar sein, und Daten-Governance
koordiniert die Speicherung in den jeweils vorgeschriebenen Regionen.

\subsection{Sovereign Cloud}
Eine Sovereign Cloud schützt personenbezogene Daten, geistiges
Eigentum, Software und Finanzdaten unter staatlicher oder regionaler
Kontrolle. Sie besteht aus drei Souveränitätsebenen:
\begin{itemize}
    \item \textbf{Daten-Souveränität:} Schlüsselkontrolle und Speicherort
    \item \textbf{Betriebs-Souveränität:} Verfügbarkeit und Regionalität
    \item \textbf{Digitale Souveränität:} Zugriffs-Policies und Transparenz
\end{itemize}

Technisch wird das umgesetzt durch Verschlüsselung,
\textit{Keep-Your-Own-Key} und Confidential Computing.
Der Betrieb erfolgt entweder über öffentliche Cloud mit regionalen
Rechenzentren (z.\,B. Frankfurt) oder über Hybrid-Plattformen wie
OpenShift. Compliance wird durch \textit{Policy-as-Code} und
kontinuierliche Überwachung sichergestellt.\\[4pt]
\textbf{Risikobasierter Ansatz:} Strenge Regeln gelten nur für
kritische Workloads. Nicht-sensible Daten dürfen EU-weit fließen,
kritische Daten (Gesundheit, Behörden, Finanzen) müssen im jeweiligen
Land bleiben, nachweisbar.

\subsubsection{Warum Sovereign Cloud? – Rechtlicher Hintergrund}

\textbf{DSGVO (Datenschutz-Grundverordnung):}
Die DSGVO ist eine EU-Verordnung (seit 2018), die den Umgang mit
personenbezogenen Daten regelt. Sie gilt für alle Unternehmen, die
Daten von EU-Bürgern verarbeiten, unabhängig davon, wo das Unternehmen
selbst sitzt. Kernanforderungen sind Zweckbindung, Datensparsamkeit,
nachweisbare Einwilligung und das Recht auf Löschung. Ergänzt wird
sie in Deutschland durch das \textbf{BDSG} (Bundesdatenschutzgesetz).

\textbf{US CLOUD Act (Clarifying Lawful Overseas Use of Data Act):}
Ein US-amerikanisches Gesetz (seit 2018), das US-Behörden erlaubt,
von amerikanischen Cloud-Anbietern (z.\,B. AWS, Azure, Google Cloud)
Herausgabe von Daten zu verlangen, auch wenn diese Daten physisch
in der EU liegen. Das untergräbt den DSGVO-Schutz, denn selbst ein
deutsches Rechenzentrum von Microsoft unterliegt dem CLOUD Act,
solange Microsoft ein US-Unternehmen ist.

\textbf{FISA 702 (Foreign Intelligence Surveillance Act):}
Ermöglicht US-Geheimdiensten den Zugriff auf Kommunikationsdaten
ausländischer Personen, die über US-amerikanische Dienste laufen,
ohne richterliche Einzelgenehmigung.\\[4pt]
Die Sovereign Cloud soll genau diese Zugriffsmöglichkeiten
durch technische und organisatorische Maßnahmen verhindern,
zum Beispiel durch europäische Betreiber ohne US-Mutterkonzern
oder durch Verschlüsselung, bei der nur der Kunde den Schlüssel hält.

\subsection{Governance}
Governance beschreibt die Möglichkeit eines Unternehmens, Anwendungen unternehmensweit 
zu steuern, Compliance nachzuweisen und Risiken zu kontrollieren. Hierfür existieren 
vordefinierte Prinzipien, Gremien und Reviews. Governance besteht aus vier Domänen:
\begin{itemize}
    \item \textbf{Geschäftsarchitektur}
    \item \textbf{Anwendungsarchitektur}
    \item \textbf{Datenarchitektur}
    \item \textbf{Technologiearchitektur}
\end{itemize}

\subsection{Personenbezogene Daten}
Daten, die eine natürliche Person identifizieren oder identifizierbar machen, sind personenbezogene Daten. Besonders schützenswert (besondere Kategorien nach Art. 9 DSGVO) sind Rasse, Ethnie, politische Meinung, Religion, Weltanschauung, Gewerkschaftszugehörigkeit, sexuelle Identität, Gesundheit sowie genetische und biometrische Daten. Daten bei unter 16-Jährigen sind nur mit der Zustimmung der Eltern zu speichern. Personenbezogene Daten können anonymisiert werden.